\bookchapter{פרק 62: פיתוח ממשק — מבוא ל-SwiftUI}

\chapterintro{הממשק הוא כל מה שהמשתמש רואה. SwiftUI הוא פריימוורק ממשק הצהרתי שאפל השיקה ב-2019 — אתם רק “מצהירים” איך הממשק אמור להיראות, והמערכת דואגת בעצמה לרינדור, להתאמה למצב הכהה ולהתאמה לכל גדלי המסך. נכון ל-2026, SwiftUI הוא כבר הבחירה המובילה באופן חד-משמעי לפרויקטים חדשים. בפרק הזה נלמד “לצייר”.}

\booksection{תפיסת העולם של SwiftUI: הצהרתי, לא פקודתי}

פיתוח ממשק ישן הוא “פקודתי”: מורים למערכת “צור כפתור שים ב-(20,30) קבע צבע”, שלב אחר שלב. SwiftUI הוא “הצהרתי”: אתם פשוט מתארים “בממשק יש כפתור אדום עם הכיתוב לחצו עליי”, ואיך לצייר ומתי לצייר מחדש — זה כבר כאב הראש של המערכת.

\begin{lstlisting}[language=Swift]
import SwiftUI

struct ContentView: View {
    var body: some View {      // body: מתאר "איך הממשק נראה"
        VStack {               // מסודרים אנכית
            Text("שלום, עולם!")
                .font(.largeTitle)      // מודיפאייר: כותרת גדולה
            Button("לחץ כאן") {
                print("נלחץ!")
            }
            .buttonStyle(.borderedProminent)  // כפתור כחול מלא
        }
        .padding()
    }
}
\end{lstlisting}

הנקודה המרכזית: תצוגה = מבנה שמקיים את פרוטוקול View; המאפיין body מתאר את הממשק; והמשנים מתשרשרים כמו שכבות “איפור”. זהו כל הבסיס של SwiftUI.

\booksection{תצוגות בסיס: טקסט, כפתור, שדה קלט ותמונה}

\begin{center}
\begin{tabularx}{\linewidth}{|X|X|X|}
\hline
\textbf{תצוגה} & \textbf{למה זה משמש} & \textbf{דוגמה} \tabularnewline
\hline
Text & הצגת טקסט & Text("שלום") \tabularnewline
Button & כפתור & Button("לחצו עליי") \{ פעולה \} \tabularnewline
TextField & שדה קלט & TextField("רמז", text: \$טקסט) \tabularnewline
Image & הצגת תמונה & \texttt{Image(systemName: "star.\allowbreak fill")} \tabularnewline
\hline
\end{tabularx}
\end{center}

\begin{center}\small\color{gray}ארבע תצוגות הבסיס של SwiftUI\end{center}

שימו לב ל-Image(systemName:) — אפל משלבת מובנה את ספריית האייקונים SF Symbols (אלפי אייקונים חינמיים), ושם כמו “star.fill” מציג ישר כוכב מלא. זו הטבה נסתרת של פיתוח iOS: כמעט ואין צורך לחפש חומרי אייקונים לממשק.

\booksection{שלישיית הפריסה: VStack, HStack, ZStack}

\begin{itemize}
  \item VStack: מסדרים אנכית (כמו תור למזנון);
  \item HStack: מסדרים אופקית (כמו ישיבה זה לצד זה);
  \item ZStack: עורמים זה על זה (כמו חפיסת קלפים; מתאים לתגיות פינה ולכיסויים).
\end{itemize}

\begin{lstlisting}[language=Swift]
VStack(alignment: .leading, spacing: 12) {
    HStack {
        Text("חשבון")
        Spacer()                 // מרווח אלסטי: דוחף את הצד הימני עד הסוף
        Text("-25.5 יואן").foregroundColor(.red)
    }
    Divider()                    // קו הפרדה
    Text("הערה: ארוחת צהריים")
}
.padding(16)                     // ריווח מכל הצדדים
.background(Color(.systemGray6)) // רקע אפור בהיר
.cornerRadius(12)                // פינות מעוגלות

// סדר המודיפאיירים משנה: קודם padding ואז background; הרקע לא כולל את השוליים
\end{lstlisting}

\begin{tcbwarn}{שימו לב: סדר המשנים משנה}
ל-padding(16).background(אפור) ול-background(אפור).padding(16) השפעות שונות — הראשון “קודם ריווח ואז צבע” (הרקע לא כולל את השוליים), והשני “קודם צבע ואז ריווח” (הרקע כולל את השוליים). מכאן נובעים “השוליים הלבנים המוזרים” הכי שכיחים אצל מתחילים. כשלא בטוחים — החליפו סדר, והתצוגה המקדימה תענה מיד.

\end{tcbwarn}

\booksection{מצב וקישור: נותנים לממשק “חיים”}

ממשק סטטי זה משעמם. הכלל של SwiftUI: הממשק תמיד משקף נאמנה את המצב; המצב השתנה — הממשק מצויר מחדש אוטומטית.

\begin{lstlisting}[language=Swift]
struct CounterView: View {
    @State private var count = 0   // @State: אם הערך הזה משתנה צריך לצייר מחדש

    var body: some View {
        VStack {
            Text("לחצת \(count) פעמים")
            Button("לחץ +1") {
                count += 1          // משנים את המצב והממשק מתרענן אוטומטית
            }
        }
    }
}
\end{lstlisting}

שדה קלט דורש “קישור דו-כיווני” — הסימן \$ מעביר פנימה גם את הערך וגם את “זכות השינוי”:

\begin{lstlisting}[language=Swift]
struct InputView: View {
    @State private var text = ""

    var body: some View {
        TextField("הקלד משהו", text: $text)  // קישור דו-כיווני עם $
            .textFieldStyle(.roundedBorder)
        Text("הקלדת: \(text)")   // מתעדכן בזמן אמת בזמן ההקלדה
    }
}
\end{lstlisting}

פעולת המשתמש

המצב משתנה (count += 1)

‏SwiftUI מזהה את השינוי

מצייר מחדש (body מחושב מחדש)

\begin{center}\small\color{gray}איור 3: זרימת הנתונים החד-כיוונית “מצב ממשק” של SwiftUI\end{center}

\booksection{תצוגה מקדימה, ערכות עיצוב ואנימציה: “הצ׳יטים” של SwiftUI}

\begin{itemize}
  \item מאקרו \#Preview: מוסיפים בתחתית הקובץ \#Preview \{ ContentView() \}, והקנבס בצד ימין מציג בזמן אמת — משנים קוד ורואים מיד;
  \item מצב כהה בחינם: משתמשים בצבעי המערכת (Color(.systemGray6), .primary) במקום לקבע צבעים, וההתאמה למצב הכהה מתבצעת אוטומטית בלי שורת קוד נוספת;
  \item אנימציה בשורה אחת: withAnimation \{ count += 1 \} או הוספת .animation(.default, value: count) לתצוגה — מעברים חלקים כמשי.
\end{itemize}

\begin{lstlisting}[language=Swift]
struct AnimatedHeartView: View {
    @State private var liked = false

    var body: some View {
        Button {
            withAnimation(.spring(duration: 0.4)) {
                liked.toggle()
            }
        } label: {
            Image(systemName: liked ? "heart.fill" : "heart")
                .font(.system(size: 48))
                .foregroundColor(liked ? .red : .gray)
                .scaleEffect(liked ? 1.3 : 1.0)
        }
    }
}
\end{lstlisting}

\begin{tcbinfo}{טיפ: קודם פונקציונליות, אחר כך אנימציה}
הטעות הנפוצה של מתחילים היא “הממשק עוד לא בנוי נכון, וכבר מתלבטים על אנימציות”. קודם דאגו שהפונקציונליות עובדת, ורק בסוף הוסיפו את האפקטים — אנימציה היא קישוט על העוגה, פונקציונליות היא העוגה עצמה. זהו גם העיקרון הקבוע שלנו: “פשטות לפני הכל”.

\end{tcbinfo}

\begin{tcbpractice}{הידיים על ההגה: בונים "כפתור מצב רוח"}
\booksection{מדריך מהיר למשנים נפוצים: עמוד אחד ביד — והממשק נפתר}

\begin{center}
\begin{tabularx}{\linewidth}{|X|X|}
\hline
\textbf{מה רוצים לעשות} & \textbf{משנה} \tabularnewline
\hline
ריווח & \texttt{.\allowbreak padding(16)} \tabularnewline
פינות מעוגלות & \texttt{.\allowbreak cornerRadius(12)} \tabularnewline
צבע רקע & \texttt{.\allowbreak background(Color(.\allowbreak systemGray6))} \tabularnewline
גופן וגודל גופן & \texttt{.\allowbreak font(.\allowbreak title.\allowbreak bold())}/\texttt{.\allowbreak font(.\allowbreak system(size: 20))} \tabularnewline
צבע חזית (טקסט/אייקון) & \texttt{.\allowbreak foregroundColor(.\allowbreak red)} \tabularnewline
מילוי כל הרוחב & \texttt{.\allowbreak frame(maxWidth: .\allowbreak infinity)} \tabularnewline
גודל קבוע & \texttt{.\allowbreak frame(width: 100, height: 44)} \tabularnewline
יישור & \texttt{VStack(alignment: .\allowbreak leading)} \tabularnewline
שקיפות & \texttt{.\allowbreak opacity(0.\allowbreak 6)} \tabularnewline
צל & \texttt{.\allowbreak shadow(radius: 5)} \tabularnewline
מחוות הקשה & \texttt{.\allowbreak onTapGesture \{ .\allowbreak .\allowbreak .\allowbreak  \}} \tabularnewline
סוג מקלדת & \texttt{.\allowbreak keyboardType(.\allowbreak decimalPad)} \tabularnewline
\hline
\end{tabularx}
\end{center}

\begin{center}\small\color{gray}מדריך מהיר למשנים הנפוצים ב-SwiftUI\end{center}

לא זוכרים? לא נורא — העמוד הזה בדיוק בשביל לגלול אליו. בפועל משתמשים הכי הרבה בארבעת הראשונים: padding, פינות מעוגלות, רקע וגופן.

הפכו את פרויקט התבנית למחליף מצבי רוח: כפתור אחד שבכל לחיצה עובר במעגל בין שמח עצוב כועס, והטקסט וצבע הרקע מתחלפים יחד איתו. רמז: אחסנו במערך שלוש קבוצות (טקסט + צבע), אחסנו את האינדקס ב-@State, ובלחיצה חשבו (index + 1) \% 3. כשתסיימו, שלחו לחברים לנסות — ותקבלו את הגל הראשון של “וואו, אתה כבר בונה אפליקציות”.

\end{tcbpractice}

\begin{tcbsummary}{סיכום הפרק}
\begin{itemize}
  \item SwiftUI הוא הצהרתי: מתארים "איך זה אמור להיראות", והמערכת דואגת לצייר ולצייר מחדש
  \item תצוגה = מבנה שמקיים את פרוטוקול View; body מתאר את הממשק; משנים מתשרשרים כמו שכבות "איפור"
  \item ארבע תצוגות הבסיס: Text / Button / TextField / Image; ספריית האייקונים SF Symbols חינמית ובשפע
  \item שלישיית הפריסה: VStack אנכי, HStack אופקי, ZStack עורם; Spacer הוא מרווח גמיש
  \item @State מעיר את הממשק לחיים; \$ מספק קישור דו-כיווני לשדה הקלט
  \item \#Preview נותן "מה שרואים זה מה שמקבלים"; צבעי המערכת מתאימים עצמם למצב הכהה; withAnimation אפקט בשורה אחת
\end{itemize}

\end{tcbsummary}
